\documentclass[a4paper, 11pt]{article}
\usepackage{fontspec}
\usepackage[margin=2cm]{geometry}
\usepackage[skip=8pt]{parskip}
\usepackage{fancyhdr}
\usepackage{booktabs}


% Custom Title Information
 

\setmainfont{Georgia}


\begin{document}

\title{Mathematics by Enquiry: Project Plan}
\author{Ragav Srinivasan}
\maketitle{}

% TODO: Project Title could be more compact?
\textit{Project Title:} Nonlinear Dynamics in Computer Graphics: An exploration of Reaction Diffusion Systems in Computer Graphics.

\textit{Overview:} The current scope of the project is to study mathematical methods applicable in Computer Graphics, particularly in procedural generation and identify their shortcomings. Then introduce an alternate method using Reaction Diffusion Systems and explain the benefits and shortcomings of this method against conventional approaches. However, As I'm just learning Reaction Diffusion systems, I am planning to broaden the scope of this project. The final scope of this project will be locked in by the end of the Midsem break.

\textit{What I'll do:} An introduction to mathematical methods that are commonly applied in Computer Graphics, examine their benefits and shortcomings. Introduce an alternative method using \textit{Reaction Diffusion Systems} borrowed from Chemistry and Biology. Explain Reaction Diffusion Systems, Application to procedural generation and study what makes it better or worse compared to conventional methods.

\textit{Original Contribution:} Hopefully, showcase a simulation by Applying Reaction Diffusion Systems.


\textit{Required Reading.} Reaction Diffusion Systems, Procedural Generation

% TODO: Where can you study these topics? Are there any papers that might be relevant?
\textit{Sources.}

\begin{itemize}
	\item  Original paper on Reaction Diffusion Systems: \textit{Turing (1952), The Chemical Basis of Morphogenesis}
	\item Turk (1991), Generating Textures on Arbitrary Surfaces Using Reaction-Diffusion
\end{itemize}


\textit{Timeline} (8 hours per week, B1 and B2 represent the break)
\\
W8 [8h]: Draft Project Plan \\
B1 [8h]: Review feedback for project plan, Required Reading, Literature Review and Decide on the  Scope of the project \\
B2 [8h]: Write Introductory text to the project \\
W9 [8h]: Submit draft for feedback, Act on feedback, Research and write more sections to the project \\
W10 [8h]: Proofread previous sections, correct any mistakes or inconsistencies found in previous sections, research and add new section to the project \\
W11 [8h]: Complete final sections of the project \\
W12 [8h]: Submit project for feedback and act on feedback \\
W13 [8h]: Proofread, Final touches and polishing \\
\\
The project will be worked on regularly throughout the rest of the semester. Each week, I'm planning to devote 8 to 10 hours of my time into my project. In the initial weeks, More time will be allocated to research and literature review (6 hours) than writing (2 hours). This will even out in the middle of the timeline and in the end I expect most of the time will be spent in writing.

\pagebreak
\textit{Plan Revisions:} The project plan is expected to be revised at frequent intervals. Updates will be much more frequent during early stages of research and writing as I'm scoping ideas for my project, however once I've been able to get a good idea of what my project will be like, the project plan will be solidified and see less updates.

\end{document}













